\newcommand*{\version}{3.31}

\documentclass[12pt,
              article,
              oneside
              ]{memoir}

\usepackage{iftex}

\ifluatex
  \usepackage{fontspec}
\fi

\usepackage[
  medievallatin,% <- principal
  % english,
  % spanish,% <- principal
  % es-minimal,
  % es-noshorthands,
  % es-tabla,
]{babel}

\ifluatex
% LuaLaTeX --------------------
\defaultfontfeatures{Ligatures=TeX, Scale=MatchLowercase}
\usepackage[MainFeatures={medium,
  Ligatures=Common,
  Numbers=OldStyle,
  SmallCapsFeatures = {LetterSpace=5.0}}]
{junicodevf}
%
% \setsansfont{Libertinus Sans}
% (without the ligatures of TeX, so that -- stays two hyphens)
\setmonofont[Scale=MatchLowercase, Ligatures=TeXOff]{DejaVu Sans Mono}
% \setmonofont[Scale=MatchLowercase]{Libertinus Mono}
% \setmonofont[Scale=MatchLowercase}]{Fira Code}
\else
% PDFLaTeX ---------------------
\usepackage[T1]{fontenc}
\usepackage[osf]{libertinus}
\usepackage[scaled=0.9]{FiraMono}
% \usepackage[scaled=0.9]{cascadia-code}
% \usepackage[scaled=0.9]{roboto}
\input glyphtounicode
\pdfgentounicode=1
\fi

\usepackage[babel=true, verbose=false, tracking=true, expansion=true, protrusion=true, final, draft=false]{microtype}
\SetTracking{ encoding = *, shape = sc }{ 30 }
\SetTracking[context = notracking, ]{encoding = *}{0}

\usepackage[autostyle=false, style=british]{csquotes}
\usepackage{latexcolors}

\newcommand{\Rojo}{\color[HTML]{8B0000}}
\newcommand*{\rojo}[1]{\textcolor[HTML]{8B0000}{#1}}
\newcommand*{\rojoit}[1]{\textit{\textcolor[HTML]{8B0000}{#1}}}
\newcommand{\Azul}{\color{bluenivaca}}
\newcommand*{\azul}[1]{\textcolor{bluenivaca}{#1}}
\newcommand{\Verde}{\color{darkgreen}}
\newcommand*{\verde}[1]{\textcolor{darkgreen}{#1}}
\newcommand{\Amarillo}{\color{goldenpoppy}}
\newcommand*{\amarillo}[1]{\textcolor{goldenpoppy}{#1}}
\newcommand{\Morado}{\color{byzantium}}
\newcommand*{\morado}[1]{\textcolor{byzantium}{#1}}
\newcommand{\Naranja}{\color{orange}}
\newcommand*{\naranja}[1]{\textcolor{orange}{#1}}

% \usepackage[margin=3cm]{geometry}

\usepackage{graphicx}
\usepackage{booktabs}

\clubpenalty=9999
\widowpenalty=9999
\brokenpenalty=5000
\frenchspacing
\raggedbottom
\midsloppy
\parindent=14pt
\fboxsep=2pt

\renewcommand\thesection{\arabic{section}}
\chapterstyle{reparticle}

\usepackage{enumitem}
\setlist{noitemsep, nosep}
\setlist[description]{font=\itshape}

\newcommand*{\elipsis}{\,.\,.\,.\,}
\newcommand*{\elipsisb}{\textup{[\elipsis]}}

\renewenvironment{quote}%
 {\begin{list}{}%
              {\setlength\rightmargin{0pt}%
              \setlength\leftmargin{\parindent}}%
  \item[]\small\ignorespaces}
 {\unskip\end{list}}


\definecolor{Link}{HTML}{408080}
\usepackage{hyperref}
\hypersetup{
  colorlinks=true,
  linkcolor=blue,
  filecolor=blue,
  citecolor = black,
  urlcolor=Link,
}
\urlstyle{tt}
\usepackage{bookmark}

\setlength{\droptitle}{-2cm}

\title{\bfseries%
  \texttt{py-words} v.\version{} \\
  Manuale usuarii}

\author{%
  Nicolás Vaughan\\
  \begin{minipage}[t]{0.8\linewidth}
    \centering\small%
    \href{mailto:nicolasvaughan@proton.me}{\texttt{nicolasvaughan@proton.me}}\\
    \url{https://nicolasvaughan.org}\\
    \url{https://github.com/nivaca/py-words}\\
    \url{https://codeberg.org/nivaca/py-words}\\
  \end{minipage}}


\date{2026-11-22}

% The pictures, made by tools/screenshots.py
\graphicspath{{img/}}
\usepackage{longtable}

% the title and the author in the PDF's own metadata
\hypersetup{pdftitle={py-words: manuale usuarii}, pdfauthor={Nicolás Vaughan}}

% the index (makeindex, with the style manual.ist: see latexmkrc)
\makeindex

% menus, keys, and the pictures of the window
\newcommand*{\menu}[1]{\textit{#1}}
\newcommand*{\then}{\,>\,}
\newcommand*{\key}[1]{\texttt{#1}}
\newcommand*{\shot}[3]{%
  \begin{figure}[htbp]
    \centering
    \includegraphics[width=#2\textwidth,height=0.75\textheight,keepaspectratio]{#1}
    \caption{#3}
    \label{fig:#1}
  \end{figure}}
% the tilde of the interactive commands (~E, ~L ...)
\newcommand*{\tl}{\textasciitilde}

\begin{document}
\maketitle

\tableofcontents*

\section{Prohemium}
\noindent
\texttt{py-words} est dictionarium Latinum quod grammaticam Latinam novit.
Verbum Latinum flexum, vel totam lineam Latinam, resolvit in articulos
dictionarii et in formas quas esse potest, cum significationibus suis
Anglice vel Hispanice; et invenit quid Latine dicatur verbum Anglicum vel
Hispanicum. Quattuor modis adhiberi potest: in fenestra, in linea
mandatorum, in terminali per totam aream visualem, et in navigatro
interretiali. Est etiam in interreti, nullo instituto, ad
\url{https://nivaca.pythonanywhere.com/}.

\index{WORDS (Whitaker)}\index{Lewis et Short}\index{De Miguel}\index{Echave-Sustaeta}\index{Latham}\index{Collatinus}%
Ipsum scripsi anno 2026. Initium eius fuit WORDS, programma Guillelmi
Whitaker; sed codex et data denuo sunt composita, et dictionarium crevit:
nunc habet verba rara et definitiones Anglicas libri Lewis et Short
\textit{A Latin Dictionary}, verba, nomina propria et definitiones
Hispanicas libri R.~de~Miguel \textit{Nuevo diccionario latino-español
etimológico}, vocabularium fundamentale J.~Echave-Sustaeta, indicem
verborum mediei evi R.~E. Latham, et vocales longas Collatini. Quilibet
articulus dicit unde veniat (capitulum~\ref{sec:codes}).

\index{consuetudines huius libri}\index{macos et windows@macOS et Windows}%
Hoc manuale primo docet quomodo programma instituatur et aperiatur;
deinde fenestram, quomodo queratur et quomodo inventa legantur; postea
cetera tria programmata, optiones et editores dictionarii. Programma
loquitur Latine, Anglice vel Hispanice; hic ponuntur nomina tabellarum
et bullarum interfaciei Latine, litteris obliquis, et tituli tabelle post
signum \enquote{>}, ut \menu{Dictionarium\then Optiones\ldots}; claves
autem, sicut in claviatura scripte sunt: \key{Ctrl+F}. In macOS,
\key{Ctrl} est clavis Command. Imagines ostendunt programmata cum
interfacie Latina et significationibus Anglicis; si interfacies est
Anglica vel Hispanica, omnes tabelle, nuntii et auxilia sunt in illa
lingua (capitulum~\ref{sec:languages}).

\index{licentia}%
\texttt{py-words} est programma liberum, sub licentia GNU General Public
License, versione 3 vel posteriore. Non utitur nisi Python et eius
bibliotheca communi (cum Tk pro fenestris et curses pro terminali), unde
nichil aliud requirit; integer catalogus programmatum quibus utitur, cum
licentiis suis, est in \texttt{docs/LICENCES.md}. WORDS et eius data
Whitaker concessit \enquote{ad omnem usum}, et textus Lewis et Short venit
ex Perseus Digital Library, sub licentia CC BY-SA 4.0.

\section{De institutione}
\label{sec:install}

Quelibet versio publicatur in GitHub et in Codeberg (loca sunt in prima
pagina), cum his fasciculis:

\index{appimage@AppImage}\index{fedora (copr, rpm)@Fedora (Copr, RPM)}\index{debian et ubuntu (deb)@Debian et Ubuntu (.deb)}\index{zip}%
\begin{itemize}
\item \texttt{py-words-v.VERSIO-windows-x64-setup.exe},
  \texttt{py-words-v.VERSIO-macos-arm64.pkg} (Apple silicon),
  \texttt{py-words-v.VERSIO-macos-x86\_64.pkg} (Intel) et
  \texttt{py-words-v.VERSIO-linux-x86\_64.run}: institutores, pro
  Windows, macOS et Linux;
\item \texttt{py-words-v.VERSIO-x86\_64.AppImage}: AppImage, pro
  quolibet Linux;
\item \texttt{py-words-VERSIO-1.fc45.noarch.rpm}: fasciculus Fedore,
  signatus, cum clave que ipsum signat, \texttt{RPM-GPG-KEY-py-words};
\item \texttt{py-words\_VERSIO-1\_all.deb}: fasciculus pro Debian et
  Ubuntu, cum subscriptione sua, \texttt{py-words\_VERSIO-1\_all.deb.asc};
\item \texttt{py-words-v.VERSIO.zip}: ipsum programma, pro Linux, macOS
  et Windows.
\item \texttt{SHA256SUMS}, cum \texttt{SHA256SUMS.asc}: summae SHA-256
  fasciculorum, et subscriptio earum. Unusquisque fasciculus etiam suam
  subscriptionem iuxta se habet (\texttt{FASCICULUS.asc}), eadem clave
  factam, sicut et nota versionis; clave importata (\texttt{gpg --import
  RPM-GPG-KEY-py-words}), \texttt{gpg --verify FASCICULUS.asc FASCICULUS}
  fasciculum probat.
\end{itemize}

\subsection{Quid requiratur}

\index{python@Python}\index{tk@Tk}%
Solum Python 3.13 vel posterius. Fenestre (dictionarium et eius editor)
requirunt Tk, quod venit cum Python in Windows et cum institutoribus
python.org in macOS. In Linux autem, et cum Python ipsius Homebrew in
macOS, potest esse fasciculus separatus:

\begin{verbatim}
$ sudo dnf install python3-tkinter      (Fedora)
$ sudo apt install python3-tk           (Debian, Ubuntu)
$ brew install python-tk                (macOS, Homebrew)
\end{verbatim}

\noindent
\texttt{python3 --version} dicit quod Python habeatur; si nimis antiquum
est, programma hoc dicit et cessat. Mandata huius libri utuntur
\texttt{python3}; in Windows scribitur \texttt{python}. Institutores,
AppImage et fasciculi suum proprium Python ferunt vel adhibent, et ideo
nullo istorum egent.

\subsection{De institutoribus (Windows, macOS, Linux)}
\label{sec:installers}

\index{institutores}\index{macos et windows@macOS et Windows}%
Via simplicissima. Quilibet institutor proprium Python 3.13, cum Tk, secum
fert, et instituit programma, locos in tabella, mandatum
\texttt{py-words} (cum eodem primo argumento quo AppImage, infra) et modum
omnia removendi. Versio recentior priori substituitur; optiones, historia
et cella servantur. Sicut in fasciculis, dictionarium institutum emendari
non potest (sectio~\ref{sec:editing} dicit quomodo tamen fiat).

\begin{description}
\item[Windows] (10 vel posterius, 64 bit). Curre \texttt{.exe}.
  Instituit pro te solo, sine iuribus administratoris, nisi in prima
  pagina eius eligas instituere pro omnibus usoribus. Loquitur Anglice vel
  Hispanice. Addit fenestram tabelle Initii et, si vis, lineam mandatorum,
  versionem terminalis in pleno visu et paginam interretialem, nexum in
  mensa et mandatum \texttt{py-words} in PATH. Removetur per ``Installed
  apps'' in Optionibus, et tunc interrogat utrum etiam optiones tuas
  removeat. Non est signatus, ideo SmartScreen de eo monere potest:
  ``More info'', deinde ``Run anyway''.
\item[macOS] (11 vel posterius). Aperi \texttt{.pkg} computatri tui
  (arm64 pro Apple silicon, x86\_64 pro Intel). Ponit
  \texttt{py-words.app} in Applications et, nisi ea omittas sub
  ``Customize'', \texttt{py-words (CLI)}, \texttt{(TUI)} et
  \texttt{(Web)}, que in Terminal aperiuntur, et mandatum
  \texttt{py-words}. Non est signatus ab Apple, ideo macOS ipsum prima
  vice recusat: in System Settings, ``Privacy \& Security'', elige ``Open
  Anyway''. \texttt{sudo py-words uninstall} ipsum removet
  (\texttt{--purge}: etiam optiones, historiam et cellam tuam).
\item[Linux] (x86-64). Curre \texttt{sh py-words-v.VERSIO-linux-x86\_64.run}.
  Ab usore communi cursus, instituit pro te solo, in
  \texttt{\textasciitilde/.local/share/py-words}, cum mandato in
  \texttt{\textasciitilde/.local/bin}; cum \texttt{sudo}, pro omnibus, in
  \texttt{/opt/py-words} et \texttt{/usr/local/bin}. Addit quattuor locos
  in tabella et icones. Primo dicit quid facturus sit et interrogat:
  \texttt{--yes} nichil interrogat, \texttt{--only-gui} addit solum locum
  fenestre, \texttt{--help} ceteras optiones ostendit. \texttt{py-words
  uninstall} ipsum removet.
\end{description}

\subsection{ZIP, vel exemplar repositorii}

\index{zip}\index{git}%
ZIP explicatur ubicumque placet, vel repositorium clonatur:

\begin{verbatim}
$ git clone https://github.com/nivaca/py-words.git
\end{verbatim}

\noindent
Capsula est totum programma: nichil est construendum vel instituendum,
et operatur ubi est.

\begin{verbatim}
$ cd py-words
$ python3 words.py amicus          (prima probatio)
$ python3 words_gui.py             (fenestra)
\end{verbatim}

\subsection{De AppImage (Linux)}

\index{appimage@AppImage}%
AppImage est unus fasciculus qui suum proprium Python fert, cum Tk. Fit
exsecutabilis et aperitur:

\begin{verbatim}
$ chmod +x py-words-v.3.31-x86_64.AppImage
$ ./py-words-v.3.31-x86_64.AppImage              (fenestra)
$ ./py-words-v.3.31-x86_64.AppImage cli amicus   (in terminali)
$ ./py-words-v.3.31-x86_64.AppImage help         (cetera)
\end{verbatim}

\noindent\index{editores dictionarii!data que scribi non possunt}%
Primum argumentum eligit programma: \texttt{gui} (nisi aliud dicatur),
\texttt{cli}, \texttt{tui}, \texttt{web}, \texttt{edit},
\texttt{edit-gui}, \texttt{check} vel \texttt{convert}; quod sequitur
datur illi programmati. In AppImage scribi non potest, et ideo in ea
dictionarium emendari non potest (capitulum~\ref{sec:editing} docet
quomodo nichilominus fiat). Instrumenta ut AppImageLauncher vel Gear
Lever eam addunt tabelle systematis.

\subsection{Fedora}

\index{fedora (copr, rpm)@Fedora (Copr, RPM)}%
Via facillima est repositorium Copr. Fasciculus utitur Python et Tk
ipsius Fedore, et addit quattuor titulos tabelle systematis, cum icone:

\begin{verbatim}
$ sudo dnf copr enable nivaca/py-words
$ sudo dnf install py-words
\end{verbatim}

\noindent
Renovationes postea veniunt cum \texttt{sudo dnf upgrade}. RPM cuiuslibet
editionis etiam manu institui potest, clave eius semel importata per
\texttt{sudo rpm --import RPM-GPG-KEY-py-words}. Operatur in qualibet
Fedora cum Python 3.13 vel posteriore, quicumque sit numerus
\texttt{fc} in nomine eius. Mandatum \texttt{py-words} deinde aperit
programmata, cum eodem primo argumento sicut AppImage:

\begin{verbatim}
$ py-words                 (fenestra)
$ py-words cli amicus      (linea mandatorum)
$ py-words help            (cetera programmata)
\end{verbatim}

\noindent\index{editores dictionarii!data que scribi non possunt}%
Programma instituitur in \texttt{/usr/share/py-words}, ubi scribi non
potest; unde nec ibi dictionarium emendari potest, sicut in AppImage.

\subsection{Debian et Ubuntu}

\index{debian et ubuntu (deb)@Debian et Ubuntu (.deb)}%
Fasciculus \texttt{.deb} eadem instituit que RPM, et etiam utitur Python
et Tk systematis. Requirit Python 3.13 vel posteriorem (Debian 13, Ubuntu
25.04 vel posteriores); \texttt{apt} simul instituit \texttt{python3-tk}:

\begin{verbatim}
$ sudo apt install ./py-words_3.31-1_all.deb
\end{verbatim}

\noindent
Subscriptio eius seorsum venit, in fasciculo \texttt{.deb.asc}, facta
eadem clave qua RPM; \texttt{gpg --verify} eam probat clave importata.
Mandatum \texttt{py-words} deinde aperit programmata, sicut cum RPM, et
\texttt{sudo apt remove py-words} ipsum removet. Nullum est repositorium:
quelibet nova versio eodem modo manu instituitur.

\subsection{De loco in tabella applicationum}

\index{tituli tabelle systematis}\index{macos et windows@macOS et Windows}%
Cum ZIP vel exemplari repositorii, semel exsequere

\begin{verbatim}
$ python3 words_install.py
\end{verbatim}

\noindent
ut programma addatur tabelle systematis, cum quattuor titulis: fenestra,
linea mandatorum et versio terminalis per totam aream visualem (he due
aperiuntur in terminali), et versio interretialis (eius servitrum currit
in terminali, et navigatrum aperitur). In Linux eunt in tabellam
applicationum, in Windows in tabellam Start, in macOS in
\texttt{\textasciitilde/Applications}. \texttt{--only-gui} addit solam
fenestram, et \texttt{--uninstall} eos removet. Iura administratoris non
requirit. Tituli aperiunt programma in sua capsula; si ergo capsula
movetur, iterum exsequendum est.

\subsection{De renovatione et remotione}

\index{renovatio}\index{remotio}%
Cum git, \texttt{git pull} in capsula; sin autem, ZIP iterum ex
interreti trahitur et capsula substituitur. Si dictionarium emendatum
est, prius servetur exemplar \texttt{data/dictionary.xml}. Ut removeatur,
tolluntur tituli tabelle per \texttt{python3 words\_install.py
--uninstall} et capsula deletur; optiones, historia et cella etiam deleri
possunt (capitulum~\ref{sec:files}).

\section{De quattuor programmatibus}
\label{sec:programs}

\index{programmata}\index{fenestra}\index{linea mandatorum}\index{terminale per totam aream}\index{pagina interretialis}%
Idem dictionarium quattuor facies habet, que optiones et historiam
inquisitionum communes habent:

\begin{longtable}{@{}lp{0.62\textwidth}@{}}
\toprule
Programma & Quid sit\\
\midrule
\endhead
\texttt{words\_gui.py} & fenestra (capitulum~\ref{sec:window})\\
\texttt{words.py} & linea mandatorum, cum modo interlocutorio (capitulum~\ref{sec:cli})\\
\texttt{words\_tui.py} & fenestra, sed in terminali, per totam aream (capitulum~\ref{sec:tui})\\
\texttt{words\_web.py} & parvum servitrum, et dictionarium in navigatro (capitulum~\ref{sec:web})\\
\texttt{meanings.py} & \texttt{words.py} cum solis formis dictionarii et significationibus\\
\texttt{words\_edit\_gui.py} & editor dictionarii, in fenestra (capitulum~\ref{sec:editing})\\
\texttt{words\_edit.py} & editor dictionarii, in terminali\\
\texttt{words\_check.py} & examinat fasciculos datorum postquam manu emendati sunt\\
\texttt{words\_install.py} & addit programma tabelle systematis\\
\bottomrule
\end{longtable}

\noindent
Ex capsula programmatis aperiuntur per \texttt{python3}:

\begin{verbatim}
$ python3 words_gui.py                     (fenestra)
$ python3 words_gui.py perspicio           (que ipsum querit)
$ python3 words_gui.py -e look             (fenestra, in Latinum)
$ python3 words.py rem acu tetigisti       (querere lineam)
$ python3 words.py                         (interlocutorie)
$ python3 words_tui.py                     (in terminali)
$ python3 words_web.py                     (in navigatro)
\end{verbatim}

\noindent\index{versio in interreti}%
Versio in interreti, \url{https://nivaca.pythonanywhere.com/}, est pagina
\texttt{words\_web.py} in servitro publico. Ibi quilibet visitator suas
proprias optiones et inquisitiones priores habet, servatas in crustulis
(\textit{cookies}) sui navigatri, et nichil horum in servitro servatur.

\section{De fenestra}
\label{sec:window}

\shot{gui_la}{1}{Fenestra, in modo obscuro, \textit{perspicio} quesito.}

\index{fenestra}\index{area inquisitionis}\index{quere (bulla)}%
Descriptio~\ref{fig:gui_la} ostendit fenestram. Superius est area ubi
verba scribuntur, et bulla \menu{Quere}, que ea querit (\key{Return}
idem facit). Area servat inquisitiones priores: \key{Down}, vel eius
sagitta, eas enumerat, recentissimam primo, et una electa iterum
queritur.

\index{quere in}\index{directio inquisitionis}\index{significationes (electio)}\index{inquisitio approximata}%
Infra sunt electiones inquisitionis:

\begin{description}
\item[Quere in] \menu{Hic}, hoc dictionarium, vel \menu{Logeion},
  dictionaria Universitatis Chicagiensis (capitulum~\ref{sec:logeion}),
  cum indice dictionariorum Logeion iuxta;
\item[directio] \menu{Latine\,$\rightarrow$\,Anglice} vel
  \menu{Anglice\,$\rightarrow$\,Latine} (cum significationibus Hispanicis,
  \menu{Hispanice} loco \menu{Anglice});
\item[Significationes] lingua significationum (capitulum~\ref{sec:languages});
\item[Inquisitio approximata] cum verbum nichil invenit, inventa orthographiarum
  vicinarum (capitulum~\ref{sec:fuzzy}).
\end{description}

\noindent
He electiones servantur: proxima vice quodlibet programma incipit sicut
relictum est.

\index{inventa!nexus}\index{inventa!explanationes}\index{duplex pressus}\index{pressus dexter}%
Inventa occupant reliquam fenestram. Sunt idem textus sicut in linea
mandatorum, sed oculis dispositus: verbum quesitum maius est, forma
dictionarii in litteris crassis cum lemmate colorato, note et
annotationes in colore cinereo, et analyses formarum in ceruleo. Primum
verbum articuli est nexus qui ipsum querit, et duplex pressus in
quolibet verbo etiam ipsum querit (verbum significationum queritur in
Latino, in altera directione). Murino super notas inter uncos vel super
analysim posito (\texttt{GEN.PL}), parva capsa explicat quid significent.
Pressus dexter verbum transcribit vel querit.

\index{linea status}\index{modus obscurus}%
Linea infima dicit quid programma faciat. Ad dexteram eius, luna fenestras
obscurat (dictionarium, eius dialogos et editorem), et deinde sol eas
iterum illuminat. Iuxta eam, nota lingue interfaciei (EN, ES vel LA)
indicem linguarum aperit.

\subsection{De tabellis}

\index{tabelle}%
\begin{description}
\item[Dictionarium] \menu{Emenda dictionarium\ldots} aperit editorem
  (capitulum~\ref{sec:editing}); \menu{Optiones\ldots} ostendit omnes
  optiones (capitulum~\ref{sec:settings}); \menu{Exi} programma claudit.
\item[Emendatio] \menu{Excide}, \menu{Transcribe} et \menu{Insere}, in area
  inquisitionis vel in inventis.
\item[Aspectus] \menu{Quere in inventis\ldots}, magnitudo textus, typi
  litterarum, et \menu{Modus obscurus}.
\item[Auxilium] \menu{Auxilium}, ductus ad querendum et ad inventa legenda;
  \menu{Pagina interretialis}, pagina programmatis; \menu{De programmate}, versio et nomina
  auctorum.
\end{description}

\subsection{Quomodo textus in inventis queratur}

\index{querere in inventis}%
\menu{Aspectus\then Quere in inventis\ldots} (\key{Ctrl+F}) aperit parvam
virgam sub inventis. Quod ibi scribitur, dum scribitur, illuminatur;
\key{Return} ad locum sequentem it, \key{Shift+Return} ad priorem, et
\key{Esc} virgam claudit.

\subsection{De textu et typis litterarum}

\index{typi litterarum}\index{magnitudo textus}\index{monitoria subtilia}%
\key{Ctrl++} et \key{Ctrl+-} textum maiorem et minorem faciunt, et
\key{Ctrl+0} magnitudinem consuetam reddit. \menu{Aspectus\then Typi litterarum\ldots}
eligit typos litterarum: unum latitudinis fixe pro inventis, quia
columne eorum ordinate manere debent, et alterum pro reliqua fenestra.
Ambo servantur. In monitorio subtili fenestre sequuntur scalam
desktopii.

\section{De inquisitione}
\label{sec:searching}

\subsection{A Latino}

\index{inquisitio!a Latino}\index{littere maiuscule}\index{u et v, i et j}\index{signa quantitatis}%
Scribitur verbum Latinum, vel tota linea Latina: quodlibet verbum suo
ordine resolvitur. Nichil refert de litteris maiusculis, nec de u et v,
nec de i et j (\textit{iulius}, \textit{Julius} et \textit{IVLIVS} sunt
idem verbum), nec de signis quantitatis et aliis accentibus:
\textit{amīcus} est \textit{amicus}, et \textit{cælum} est
\textit{caelum}.

\index{artificia}\index{prefixa et suffixa}\index{enclitica}\index{orthographie mediei evi}%
Verbum primo queritur tamquam radix et terminatio dictionarii. Si sic
nichil invenitur, programma temptat prefixa, suffixa et enclitica
(\textit{-que}, \textit{-ne}, \textit{-ve}), et deinde alias
orthographias: perfectum syncopatum, duo verba tamquam unum scripta
(\textit{rempublicam}), \textit{-st} finale pro \textit{est}
(\textit{factumst}), et orthographias Latinitatis mediei evi (\textit{e}
pro \textit{ae}, \textit{y} pro \textit{i}, \textit{michi} pro
\textit{mihi}). Semper dicit quid fecerit:

\begin{verbatim}
hystorie
  spelling: Medieval hystorie -> historiae
    'hystorie' may be a medieval spelling of historiae
  historia, historiae  noun, 1st declension, feminine  [XXXAX]
    a narrative of past events, history; ...
    histori-ae     GEN.SG
\end{verbatim}

\noindent\index{nomina propria}%
Verbum ignotum cum littera maiuscula habetur pro nomine proprio, et
relinquitur.

\subsection{In Latinum}

\index{inquisitio!in Latinum}\index{partes orationis}%
In altera directione (\menu{Anglice\,$\rightarrow$\,Latine} in fenestra,
\texttt{-e} in linea mandatorum), scribitur unum verbum lingue
significationum: enumerantur articuli Latini quorum significationes ipsum
habent, optimi primo, sex ad summum, et deinde quot alii sint. Pars
orationis post verbum solam illam relinquit: \texttt{friend ADJ}. Partes
orationis sunt \texttt{N}, \texttt{PRON}, \texttt{ADJ}, \texttt{NUM},
\texttt{ADV}, \texttt{V}, \texttt{PREP}, \texttt{CONJ} et
\texttt{INTERJ}. Verba valde communia (\textit{the}, \textit{a}) numquam
queruntur, et Hispanice accentus nichil referunt: \textit{corazon}
invenit \textit{corazón}.

\subsection{De inquisitione approximata}
\label{sec:fuzzy}

\index{inquisitio approximata}\index{errores scribendi}%
Communiter verbum cum errore scribendi nichil invenit. Cum inquisitione
approximata, verbum trium litterarum vel plurium quod adhuc nichil
invenit sequuntur inventa verborum que una mutatione ab ipso distant
(una littera omissa, addita vel mutata, vel due littere permutate), tria
ad summum, primo ea articulorum frequentiorum:

\begin{verbatim}
amicuss
  (unknown)
  (near spellings: amicus)

amicus
  amicus, amici  noun, 2nd declension, masculine  [XXXBO]
  ...
\end{verbatim}

\noindent
Inquisitio in Latinum idem facit cum verbis significationum:
\textit{freind} invenit \textit{friend}, et \textit{wolfs} invenit
\textit{wolf}. Agit nisi auferatur: \menu{Inquisitio approximata} in fenestra et in
pagina interretiali, \key{F3} in terminali, vel \texttt{--no-fuzzy} in
linea mandatorum, pro una vice.

\section{Quomodo inventa legantur}
\label{sec:results}

\index{inventa}%
Hoc est quod \texttt{py-words} ostendit pro \textit{amicorum}:

\begin{verbatim}
amicorum
  amicus, amici  noun, 2nd declension, masculine  [XXXBO]
    A friend, patron, protector, companion; ... (Lewis and Short)
    amīc-ōrum      GEN.PL
  amicus, amica -um, amicior -or -us, amicissimus -a -um
      adjective, 1st/2nd declension  [XXXAO]
    friendly, kind, amicable, favorable, ... (Lewis and Short)
    amīc-ōrum      GEN.PL.M
    amīc-ōrum      GEN.PL.N
\end{verbatim}

\index{inventa!linee}\index{partes principales}\index{vocales longe}%
\begin{itemize}
\item Prima linea est verbum quesitum. Sequitur, cum indentatione,
  quilibet articulus quem esse potest.
\item Prima linea articuli habet formam dictionarii (partes principales:
  \textit{amicus, amici}), partem orationis et grammaticam verbis
  expressam (\textit{noun, 2nd declension, masculine}, id est nomen
  secunde declinationis, masculinum), et notas inter uncos
  (capitulum~\ref{sec:codes}); quidam articuli addunt etatem vel
  frequentiam verbis expressam (\textit{archaic}, \textit{lesser}\ldots).
\item Linea sequens est significatio. Eius sensus per punctum et virgulam
  separantur, communissimo primo. In fine, dictionarium unde venit, si non
  est WORDS: (Lewis and Short), (De Miguel), (Echave-Sustaeta).
  \texttt{[en]} ante significationem dicit quod non est translata, sed
  Anglice datur.
\item Deinde veniunt forme quas verbum esse potest: radix et terminatio,
  per lineolam iuncte, cum vocalibus longis notatis ubi noscuntur
  (\textit{amīc-ōrum}), et que forma sit, per abbreviationes capituli~\ref{sec:glosses}.
  Eadem scriptura potest esse plures forme: \textit{amīc-ōrum} adiectivi
  est masculinum et neutrum.
\item Linea ante articulum potest dicere quomodo verbum inventum sit:
  prefixum, suffixum vel encliticum, cum significatione sua, vel
  mutatio orthographie.
\end{itemize}

\noindent\index{tempora composita}%
Participium quod sequitur forma verbi \textit{esse} legitur tamquam
tempus compositum:

\begin{verbatim}
laudatus
  laudo, laudare, laudavi, laudatus  verb, 1st conjugation  [XXXAX]
    to praise, laud, commend, extol, ... (Lewis and Short)
    laudāt-us      PRF.PASS.PTCP.NOM.SG.M
    participle + est  3SG.PRF.PASS.IND
      PERF PASSIVE PPL + verb TO_BE => PASSIVE perfect system
\end{verbatim}

\noindent\index{inventa!annotationes}%
Quedam annotationes inter parentheses ponuntur:

\begin{longtable}{@{}>{\raggedright}p{0.4\textwidth}p{0.54\textwidth}@{}}
\toprule
Annotatio & Quid dicat\\
\midrule
\endhead
(unknown) & verbum non est inventum\\
(near spellings: \ldots) & sequuntur inventa inquisitionis approximate\\
(less likely forms are not shown) & optio \texttt{trim\_output} quasdam omisit\\
(no match) & inquisitio in Latinum nichil invenit\\
(near words: \ldots) & inquisitio approximata in Latinum\\
(\ldots\ more not shown) & post sex optimos articulos, in Latinum\\
\bottomrule
\end{longtable}

\subsection{De notis inter uncos}
\label{sec:codes}

\index{note inter uncos}\index{etas}\index{materia}\index{regio}\index{frequentia}\index{fontes articulorum}%
Quinque littere inter uncos post articulum dicunt quando, ubi, quam sepe
et qua auctoritate verbum adhibeatur, hoc ordine: etas eius, materia (res
ad quam pertinet), regio, frequentia et fons. X tantum dicit quod nichil
singulare est dicendum. Sic \texttt{[XXXBO]} vocabuli \textit{amicus}
legitur: in usu per omnes etates, generale, ubique, frequens, et sumptum
ex Oxford Latin Dictionary. Articulus potest plures fontes habere: ceteri
sequuntur quinque litteras, quilibet post signum \texttt{+}
(\texttt{[XXXBO+A]} etiam ex Echave-Sustaeta venit). Etas, si noscitur,
et frequentia, si minor vel rarior est, etiam verbis Anglicis post notas
scribuntur. Littere sunt he:

\begin{longtable}{@{}lp{0.8\textwidth}@{}}
\toprule
Littera & Etas\\
\midrule
\endhead
X & in usu per omnes etates\\
A & archaica\\
B & antiqua (ante tempus classicum)\\
C & classica\\
D & sera (post annum 200)\\
E & posterior (post annum 600)\\
F & mediei evi (post annum 1000)\\
G & doctorum (a renovatione litterarum et postea)\\
H & moderna\\
\bottomrule
\end{longtable}

\begin{longtable}{@{}lp{0.8\textwidth}@{}}
\toprule
Littera & Materia\\
\midrule
\endhead
X & generalis\\
A & agricultura, plante, animalia, terra\\
B & biologia, medicina, corpus\\
D & drama, musica, artes\\
E & Ecclesia, Biblia\\
G & grammatica, rhetorica, logica, litterature, schola\\
L & ius, regimen, pecunia, tituli\\
P & poesis\\
S & scientia, philosophia, mathematica, mensure\\
T & artes mechanice, architectura, agrimensura\\
W & bellum, exercitus, classis\\
Y & mythologia\\
\bottomrule
\end{longtable}

\begin{longtable}{@{}lp{0.8\textwidth}@{}}
\toprule
Littera & Regio\\
\midrule
\endhead
X & ubique\\
A & Africa\\
B & Britannia\\
C & Sina\\
D & Scandinavia\\
E & Egyptus\\
F & Francia, Gallia\\
G & Germania\\
H & Grecia\\
I & Italia, Roma\\
J & India\\
K & Balcani\\
N & Belgium (Terre Basse)\\
P & Persia\\
Q & Oriens Propinquus\\
R & Russia\\
S & Hispania\\
U & Europa orientalis\\
\bottomrule
\end{longtable}

\begin{longtable}{@{}lp{0.8\textwidth}@{}}
\toprule
Littera & Frequentia\\
\midrule
\endhead
X & ignota\\
A & frequentissima\\
B & frequens\\
C & communis\\
D & minor\\
E & non communis\\
F & rarissima\\
I & in inscriptionibus tantum\\
M & in graphitis tantum\\
N & apud Plinium tantum\\
\bottomrule
\end{longtable}

\begin{longtable}{@{}lp{0.8\textwidth}@{}}
\toprule
Littera & Fons\\
\midrule
\endhead
X & generalis vel ignotus\\
A & Echave-Sustaeta, \textit{Lengua latina. Vocabulario básico}\\
B & Beeson, \textit{Medieval Latin}\\
C & \textit{Cassell's Latin Dictionary}\\
D & Adams, \textit{Latin Sexual Vocabulary}\\
E & Stelten, \textit{Ecclesiastical Latin}\\
F & Deferrari, de sancto Thoma de Aquino\\
G & Gildersleeve et Lodge, \textit{Latin Grammar}\\
H & Collatinus (Ouvrard)\\
I & Leverett, \textit{Lexicon}\\
J & Bracton, \textit{De Legibus}\\
K & \textit{Calepinus Novus} (Licoppe)\\
L & Lewis, \textit{Elementary Latin Dictionary}\\
M & Latham, \textit{Medieval Word List}; DMLBS\\
N & Nelson, index verborum\\
O & \textit{Oxford Latin Dictionary}\\
P & Souter, \textit{Later Latin}\\
Q & alia dictionaria\\
R & Plater et White, \textit{Vulgate grammar}\\
S & Lewis et Short, \textit{A Latin Dictionary}\\
T & inventum in translatione\\
U & Du Cange\\
V & Blatt, Saxo\\
W & coniectura Whitaker\\
Y & ad tempus\\
Z & ab usuario missum\\
m & De Miguel, \textit{Nuevo diccionario latino-español etimológico}\\
\bottomrule
\end{longtable}

\subsection{De abbreviationibus analysium}
\label{sec:glosses}

\index{abbreviationes glossarum}\index{Leipzig Glossing Rules}%
Forme resolvuntur per abbreviationes regularum glossandi Lipsiensium
(\textit{Leipzig Glossing Rules}), punctis iunctas; persona et numerus
simul eunt, ita quod \texttt{2SG.PRF.ACT.IND} est secunda persona
singularis perfecti activi indicativi. Que signo * notantur non sunt in
catalogo Lipsiensi. Optione \texttt{use\_glosses} ablata
(\texttt{--no-glosses} in linea mandatorum), analyses verbis Anglicis
scribuntur: \enquote{accusative singular}.

\begin{longtable}{@{}lp{0.6\textwidth}@{}}
\toprule
Abbreviatio & Significatio\\
\midrule
\endhead
\texttt{NOM}, \texttt{GEN}, \texttt{DAT} & nominativus, genitivus, dativus\\
\texttt{ACC}, \texttt{ABL}, \texttt{VOC}, \texttt{LOC} & accusativus, ablativus, vocativus, locativus\\
\texttt{SG}, \texttt{PL} & singularis, pluralis\\
\texttt{M}, \texttt{F}, \texttt{N}, \texttt{C}* & masculinum, femininum, neutrum, commune\\
\texttt{1}, \texttt{2}, \texttt{3} & prima, secunda, tertia persona\\
\texttt{PRS}, \texttt{IMPF}*, \texttt{FUT} & presens, imperfectum, futurum\\
\texttt{PRF}, \texttt{PLUP}*, \texttt{FUT.PRF}* & perfectum, plusquamperfectum, futurum exactum\\
\texttt{ACT}*, \texttt{PASS} & activum, passivum\\
\texttt{IND}, \texttt{SBJV}, \texttt{IMP} & indicativus, subiunctivus, imperativus\\
\texttt{INF}, \texttt{PTCP}, \texttt{SUP}* & infinitivus, participium, supinum\\
\texttt{POS}*, \texttt{CMPR}*, \texttt{SPRL}* & positivus, comparativus, superlativus\\
\texttt{CARD}*, \texttt{ORD}* & numerale cardinale, ordinale\\
\texttt{DISTR}, \texttt{MULT}* & numerale distributivum, multiplicativum (adverbium numerale)\\
\bottomrule
\end{longtable}

\section{De linguis}
\label{sec:languages}

\subsection{De lingua significationum}

\index{significationes (electio)}\index{significationes Hispanice}\index{De Miguel}\index{Echave-Sustaeta}%
Significationes sunt Anglice vel Hispanice, et similiter cetera inventa
(\enquote{sustantivo, 2.ª declinación, masculino}); abbreviationes
glossarum in utraque eedem sunt. Significationes Hispanice veniunt ex
dictionario De Miguel ubicumque verbum habet, sicut ipse eas scripsit
anno 1867, cum sua orthographia (\textit{razon}, \textit{á}), et hoc in
fine dicunt, \enquote{(De Miguel)}; pauce ex Echave-Sustaeta veniunt;
cetere ex Anglico translate sunt. Eligitur per \menu{Significationes} in
fenestra et in pagina interretiali, per \key{F5} in terminali, vel per
\texttt{--language es} in linea mandatorum, pro una vice. Inquisitio in
Latinum incipit a lingua electa.

\subsection{De lingua interfaciei}

\index{lingua interfaciei}%
Programmata loquuntur Latine, Anglice vel Hispanice: tabelle, bulle,
nuntii, auxilium et optiones, editores et examinator. Lingua eligitur
per bullam cum nota eius (EN, ES vel LA), in parte dextra inferiore
fenestre et in summa pagina interretiali, per \menu{Lingua interfaciei}
in optionibus fenestre et pagine interretialis, per \key{F6} in
terminali (nota eius est ad dexteram linee superioris), vel in fine dialogi
\texttt{\#} linee mandatorum (\texttt{--interface la} pro una vice; nota
lingue Latine est \texttt{la}). Fenestra et pagina interretialis statim
mutantur. Mutata interfacie, significationes quoque in eandem linguam
mutantur, si dictionarium eam habet; significationes autem Latine non
sunt, ita quod cum interfacie Latina manent Anglice vel Hispanice, sicut
erant, et similiter annotationes inventorum, ut \enquote{(unknown)}.
Postea utraque separatim eligi potest.

\section{De Logeion}
\label{sec:logeion}

\index{Logeion}%
Loco articulorum huius dictionarii, programmata possunt ostendere
articulos Logeion (\url{https://logeion.uchicago.edu}), dictionaria
Latina et Greca Universitatis Chicagiensis, cum definitione brevi,
frequentia et exemplis ex corpore suo. Interrete requirit. Eligitur per
\menu{Quere in} in fenestra et in pagina interretiali, per \key{F8} in
terminali, per \texttt{\tl G} in linea mandatorum interlocutoria, vel per
\texttt{--logeion} pro una vice.

\index{lemmata}%
Logeion formas flexas non resolvit. Ideo verba hic primo resolvuntur, et a
Logeion petuntur lemmata eorum: \textit{dixit} ostendit \textit{dico}, et
\textit{est} ostendit \textit{edo} et \textit{sum}. Quodlibet verbum
primo ostendit analysim suam, sine significationibus, et deinde paginas
Logeion lemmatum suorum. Dum Logeion electum est, inquisitio a solo
Latino incipit, et altera directio, lingua significationum, inquisitio
approximata et editor colore cinereo obducuntur.

\index{Logeion!dictionaria}\index{DMLBS}%
Articuli possunt esse unius dictionarii Logeion, vel omnium; donec aliud
eligatur, est DMLBS. Eligitur per indicem iuxta \menu{Quere in}, per
\key{F5} in terminali (dum Logeion electum est), per \texttt{\tl D} in
linea mandatorum interlocutoria, vel per \texttt{--dictionary} cum nota
sua:

\begin{longtable}{@{}lp{0.7\textwidth}@{}}
\toprule
Nota & Dictionarium\\
\midrule
\endhead
\texttt{all} & omnia dictionaria\\
\texttt{ls} & Lewis \& Short\\
\texttt{lewis} & Lewis, \textit{Elementary Latin Dictionary}\\
\texttt{georges} & Georges\\
\texttt{gaffiot} & Gaffiot 2016\\
\texttt{dmlbs} & DMLBS (\textit{Dictionary of Medieval Latin from British Sources})\\
\texttt{lane} & LaNe\\
\texttt{ducange} & Du Cange\\
\texttt{lana} & Lana-Michelone\\
\texttt{frieze} & Frieze-Dennison (Vergilius)\\
\texttt{sinicum} & Latino-Sinicum\\
\bottomrule
\end{longtable}

\noindent
Mutato dictionario, inventa statim iterum ostenduntur, Logeion non
iterum interrogato. In fenestra, nexus articuli ad alia verba illa verba
querunt, et ceteri nexus in navigatro aperiuntur.

\section{De linea mandatorum}
\label{sec:cli}

\shot{cli_la}{1}{Linea mandatorum, interlocutorie, cum \textit{perspicio}.}

\index{linea mandatorum}%
\texttt{words.py} querit verba post ipsum data, tamquam unam lineam
Latinam, vel, sine verbis, modum interlocutorium incipit
(descriptio~\ref{fig:cli_la}). \texttt{meanings.py} idem est quod
\texttt{words.py -m}.

\begin{verbatim}
$ python3 words.py rem acu tetigisti
$ python3 words.py -e friend --pos N
$ python3 words.py -f textus.txt -o inventa.txt
$ python3 words.py --language es amicorum
$ python3 words.py --logeion arma
$ python3 words.py --dictionary gaffiot arma
\end{verbatim}

\index{linea mandatorum!optiones}%
\begin{longtable}{@{}>{\raggedright}p{0.36\textwidth}p{0.58\textwidth}@{}}
\toprule
Optio & Quid faciat\\
\midrule
\endhead
\texttt{-f FILE} & querit quamlibet lineam fasciculi FILE (\texttt{-} pro introitu communi)\\
\texttt{-o FILE} & scribit inventa in FILE\\
\texttt{-e} & a lingua significationum in Latinum\\
\texttt{--pos POS} & cum \texttt{-e}, solam hanc partem orationis\\
\texttt{-m} & solas formas dictionarii et significationes\\
\texttt{--no-glosses} & analyses verbis, non abbreviationibus\\
\texttt{--no-fuzzy}, \texttt{--fuzzy} & sine vel cum inquisitione approximata\\
\texttt{--logeion}, \texttt{--local} & Logeion, vel hoc dictionarium\\
\texttt{--dictionary DICT} & dictionarium Logeion (capitulum~\ref{sec:logeion})\\
\texttt{--dictionaries} & enumerat dictionaria Logeion\\
\texttt{--language CODE} & lingua significationum (\texttt{en}, \texttt{es})\\
\texttt{--languages} & enumerat linguas significationum\\
\texttt{--interface CODE} & lingua interfaciei\\
\texttt{--datadir DIR} & fasciculi datorum in capsula DIR\\
\texttt{--guide} & imprimit auxilium: quomodo inventa legantur\\
\texttt{--config} & facit fasciculum configurationis et dicit ubi sit\\
\texttt{--clear-history} & obliviscitur inquisitionum priorum\\
\texttt{--version} & versio et nomina auctorum\\
\bottomrule
\end{longtable}

\noindent
Optiones que aliquid eligunt (glossas, inquisitionem approximatam, fontem,
linguas) hoc faciunt pro illa sola vice; optiones servate retinent quod
habebant.

\subsection{De modo interlocutorio}

\index{linea mandatorum!modus interlocutorius}%
In modo interlocutorio queritur quelibet linea scripta post signum
\texttt{=>}, nisi incipiat ab aliquo istorum:

\begin{longtable}{@{}>{\raggedright}p{0.36\textwidth}p{0.58\textwidth}@{}}
\toprule
Linea & Quid faciat\\
\midrule
\endhead
\texttt{\tl E}, \texttt{\tl L} & ab Anglico in Latinum; a Latino in Anglicum\\
\texttt{\tl G} & Logeion vel hoc dictionarium\\
\texttt{\tl D [DICT]} & dictionarium Logeion (solum: index ex quo eligatur)\\
\texttt{@FILE} & querit quamlibet lineam fasciculi FILE\\
\texttt{\#} & mutat parametra (capitulum~\ref{sec:settings})\\
\texttt{!} & mutat modos elaboratoris\\
\texttt{\tl H} & auxilium: quomodo queratur et inventa legantur\\
\texttt{\tl A} & de programmate: versio et nomina auctorum\\
\texttt{?} & index harum linearum\\
\bottomrule
\end{longtable}

\noindent\index{historia inquisitionum}%
\key{Up} et \key{Down} reducunt inquisitiones priores, etiam sessionum
preteritarum et aliorum programmatum. In dialogis \texttt{\#} et
\texttt{!}, \texttt{?} quamlibet optionem explicat. Due linee vacue, vel
\key{Ctrl+D}, programma claudunt. \key{Ctrl+C} non claudit: sistit quod
fit (longum catalogum inventorum, fasciculum), et programma lineam
sequentem expectat.

\subsection{De paginatore et coloribus}

\index{paginator}\index{colores}%
In terminali, inventa longiora quam area visualis aperiuntur in
paginatore systematis, \texttt{less} (vel eo quem \texttt{\$PAGER}
nominat), sicut cum \texttt{man}: \key{PgUp} et \key{PgDn} ea movent,
\texttt{/} textum querit, et \texttt{q} ad signum \texttt{=>} redit.
Inventa sunt colorata, coloribus ipsius terminalis, sicut in fenestra.
Inventa in fasciculum vel in tubulum scripta sunt simplicia, et
\texttt{NO\_COLOR=1} colores aufert.

\section{De terminali per totam aream}
\label{sec:tui}

\shot{tui_la}{1}{\texttt{words\_tui.py}, in terminali.}

\index{terminale per totam aream}%
\texttt{words\_tui.py} est fenestra in terminali
(descriptio~\ref{fig:tui_la}): linea querendi superius, inventa infra, et
claves in ultima linea. \key{Up} et \key{Down} reducunt inquisitiones
priores, et \key{PgUp} et \key{PgDn} inventa movent. In Windows requirit
fasciculum \texttt{windows-curses}.

\index{terminale per totam aream!claves}%
\begin{longtable}{@{}>{\raggedright}p{0.3\textwidth}p{0.64\textwidth}@{}}
\toprule
Clavis & Actio\\
\midrule
\endhead
\key{F1} & auxilium, et nomina auctorum\\
\key{F2}, \key{Ctrl+L} & a Latino, vel in Latinum\\
\key{F3}, \key{Ctrl+O} & optiones (\key{Space} unam mutat)\\
\key{F4}, \key{Ctrl+E} & editor dictionarii\\
\key{F5}, \key{Ctrl+G} & lingua significationum (cum Logeion, eius dictionarium)\\
\key{F6} & lingua interfaciei\\
\key{F7} & omnia inventa transcribit\\
\key{F8} & hoc dictionarium vel Logeion\\
\key{F9} & historiam delet\\
\key{F10}, \key{Ctrl+Q} & programma claudit\\
\bottomrule
\end{longtable}

\section{De pagina interretiali}
\label{sec:web}

\shot{web_la}{0.8}{\texttt{words\_web.py}, in navigatro.}

\index{pagina interretialis}%
\texttt{words\_web.py} incipit parvum servitrum et aperit navigatrum ad
\url{http://127.0.0.1:8000/}, vel ad aliam portam si hec occupata est
(descriptio~\ref{fig:web_la}). Pagina habet aream inquisitionis,
directionem, \menu{Significationes}, \menu{Inquisitio approximata}, \menu{Quere in}, et
inventa sicut in fenestra, cum eisdem nexibus et explanationibus; infra,
inquisitiones priores, omnes optiones, auxilium et nomina auctorum. Luna
paginam obscurat, et sol iterum illuminat; donec primo adhibeantur,
pagina sequitur systema. Ante lunam, bulla rotunda cum nota lingue
interfaciei (EN, ES vel LA) indicem linguarum aperit.

\index{pagina interretialis!loca}%
Quelibet inquisitio suum proprium locum habet, ut
\url{http://127.0.0.1:8000/?q=amicus+est} (\texttt{\&lang=en} ab Anglico
in Latinum), ita quod tamquam signum libri servari potest. Servitrum
dictionarium iterum legit cum editor ipsum mutat; dictionarium ex
navigatro emendari non potest.

\index{pagina interretialis!optiones}\index{rete}%
Solum hoc computatrum paginam aperire potest, nisi \texttt{--host} aliud
dicat: cum \texttt{--host 0.0.0.0} adhiberi potest ex telephono in eodem
rete. Nulla sunt verba secreta; ideo hoc fiat tantum in rete fido.
\texttt{--port} eligit portam, \texttt{--no-browser} navigatrum non
aperit, et \key{Ctrl+C} servitrum sistit.

\section{De optionibus}
\label{sec:settings}

\index{optiones}\index{parametra}%
\menu{Dictionarium\then Optiones\ldots} ostendit omnes optiones fenestre,
cum eo quod quelibet facit (murino super optionem posito, eius explanatio
apparet); mutatio statim valet. Sunt eedem que parametra (\texttt{\#}) et
modi elaboratoris (\texttt{!}) linee mandatorum, \key{F3} in terminali, et
optiones pagine interretialis. Dialogus habet etiam \menu{Lingua
interfaciei}, \menu{Typi litterarum\ldots}, \menu{Dele historiam} et \menu{Optiones
primitive}. Optiones in linea mandatorum et in fasciculo configurationis
nomina sua Anglica retinent.

\index{parametra!quid inventa ostendant}%
\begin{longtable}{@{}>{\raggedright}p{0.36\textwidth}p{0.58\textwidth}@{}}
\toprule
Parametrum & Quid faciat, cum agit\\
\midrule
\endhead
\multicolumn{2}{@{}l}{\itshape Quid inventa ostendant}\\
\texttt{show\_dictionary\_forms} & quemlibet articulum cum partibus principalibus et grammatica\\
\texttt{show\_only\_meanings} & solos articulos et significationes, sine formis\\
\texttt{use\_glosses} & analyses abbreviationibus glossarum, non verbis\\
\texttt{show\_long\_vowels} & vocales longas notatas in formis (\textit{amīc-ōrum})\\
\texttt{show\_codes} & notas inter uncos (\texttt{[XXXBO]})\\
\texttt{explain\_codes} & etatem et frequentiam semper verbis, non solum cum insolite sunt\\
\texttt{trim\_output} & omittit inventa minus probabilia (et, in Latinum, omnia preter sex optima)\\
\texttt{trim\_medieval} & cum \texttt{trim\_output}, etiam formas mediei evi si alie invente sunt\\
\addlinespace
\multicolumn{2}{@{}l}{\itshape Quomodo verba querantur}\\
\texttt{fuzzy\_search} & orthographias vicinas verbi quod nichil invenit\\
\texttt{use\_affixes} & verbum ignotum legit tamquam notum cum prefixo vel suffixo\\
\texttt{use\_tricks} & alias orthographias, et duo verba tamquam unum scripta\\
\texttt{medieval\_spellings} & cum \texttt{use\_tricks}, etiam orthographias mediei evi\\
\texttt{read\_compounds} & participium et formam verbi \textit{esse} tamquam tempus compositum\\
\texttt{ignore\_unknown\_names} & nulla artificia in verbo ignoto cum littera maiuscula\\
\texttt{near\_entries} & pro verbo ignoto, articulos vicinos loco ubi esset\\
\addlinespace
\multicolumn{2}{@{}l}{\itshape Modi elaboratoris}\\
\texttt{show\_entry\_source} & ex quo dictionario quilibet articulus sit, et numerum eius ibi\\
\texttt{pearse\_codes} & quelibet linea inventorum incipit a nota (pro aliis programmatibus)\\
\texttt{pause\_in\_screen\_output} & inventa longa in paginatore, in linea mandatorum\\
\bottomrule
\end{longtable}

\noindent\index{trim\_output}%
Poete verbis et formis insolitis utuntur, et ideo \texttt{trim\_output}
potest rectam omittere; si versus inventa mira dat, auferatur.

\section{De dictionario emendando}
\label{sec:editing}

\index{editores dictionarii}%
Duo editores addunt, corrigunt et delent articulos dictionarii,
\texttt{data/dictionary.xml}: unus in fenestra,
\texttt{words\_edit\_gui.py}, quem etiam \menu{Dictionarium\then Emenda
dictionarium\ldots} aperit, alter in terminali, \texttt{words\_edit.py}. In
utroque:

\begin{itemize}
\item articulus queritur per formam Latinam (\textit{amicorum}), radicem,
  vel verbum significationis sue;
\item corrigitur vel deletur, vel novus fit, vacuus vel similis ei qui
  ostenditur (\menu{Novus huic similis}); campi pendent ex parte orationis, et
  note eliguntur ex indicibus qui eas explicant;
\item prospectus ostendit articulum sicut inventa ipsum ostendent, et
  quodlibet vitium quod habet; \menu{Proba verbum} ostendit quid programma
  faciat de verbo cum dictionario sicut nunc est;
\item significationes Hispanice etiam emendari possunt, iuxta Anglicas,
  lingua electa per \menu{Significationes} (\key{F7} in terminali).
\end{itemize}

\index{editores dictionarii!servare}\index{fasciculi .bak}%
Mutatio intrat dictionarium per \menu{Applica mutationes} (\key{F3} in
terminali), et fasciculus scribitur tantum per \menu{Serva} (\key{F2}).
Si relinquitur articulus cum mutationibus non applicatis, editor querit
quid de eis fiat. Prima vice qua in sessione servatur, fasciculi
originales retinentur tamquam fasciculi \texttt{.bak}.

\index{editores dictionarii!claves}%
\begin{longtable}{@{}>{\raggedright}p{0.3\textwidth}p{0.64\textwidth}@{}}
\toprule
Clavis (terminale) & Actio\\
\midrule
\endhead
\key{F3}, \key{Ctrl+A} & applicat mutationes (vel addit articulum novum)\\
\key{F4}, \key{Ctrl+N} & articulus novus\\
\key{F5}, \key{Ctrl+D} & articulus novus similis huic\\
\key{F6}, \key{Ctrl+T} & verbum probat\\
\key{F7}, \key{Ctrl+L} & lingua significationum que preter Anglicam emendatur\\
\key{F8}, \key{Ctrl+K} & articulum delet\\
\key{F2}, \key{Ctrl+S} & fasciculum servat (mutationibus prius applicatis)\\
\key{F10}, \key{Ctrl+Q} & programma claudit\\
\key{Tab} & area sequens: inquisitio, index, forma articuli\\
\key{Ctrl+F} & querere\\
\bottomrule
\end{longtable}

\noindent\index{editores dictionarii!data que scribi non possunt}\index{capsula datorum}%
In AppImage et in RPM data scribi non possunt, et ideo editores non
agunt. Ut ibi dictionarium emendetur, capsula \texttt{data} alio
transcribitur, et exemplar datur tamquam capsula datorum: in fasciculo
configurationis (\texttt{folder} in \texttt{[data]}), per
\texttt{--datadir}, vel per \texttt{\$PY\_WORDS\_DATADIR}.

\index{words\_check.py}%
Fasciculi etiam manu emendari possunt, in quolibet editore textus:
transcribitur articulus verbi similis et mutatur. Deinde \texttt{python3
words\_check.py} indicat, cum numeris linearum, que programma legere non
potest et articulos qui probabiliter non faciunt quod intendebatur.

\section{De clavibus fenestre}
\label{sec:keys}

\index{claves}%
\begin{longtable}{@{}>{\raggedright}p{0.4\textwidth}p{0.54\textwidth}@{}}
\toprule
Clavis & Actio\\
\midrule
\endhead
\key{Return} & querere\\
\key{Down} (in area inquisitionis) & inquisitiones priores\\
\key{Ctrl+F} & querere in inventis\\
\key{Return}, \key{Shift+Return} (in inventis querendo) & locus sequens et prior\\
\key{Esc} (in inventis querendo) & claudere virgam inquisitionis\\
\key{Ctrl++}, \key{Ctrl+-} & textus maior et minor\\
\key{Ctrl+0} & magnitudo consueta\\
\key{Ctrl+X}, \key{Ctrl+C}, \key{Ctrl+V} & excidere, transcribere, inserere\\
\key{F1} & auxilium\\
\key{Ctrl+Q} & claudere programma\\
\key{Esc} (in dialogo) & ipsum claudere\\
\bottomrule
\end{longtable}

\section{Ubi quelibet res servetur}
\label{sec:files}

\index{fasciculi configurationis}\index{config.toml}\index{historia inquisitionum}\index{cella}%
Optiones, typi litterarum, modus obscurus et electiones fenestre in uno
fasciculo configurationis servantur, \texttt{config.toml}. Est textus
simplex, cum annotatione pro qualibet optione, et manu emendari potest;
\texttt{python3 words.py --config} ipsum facit et dicit ubi sit.
Programma etiam servat ultimas mille inquisitiones, et cellam dictionarii
per quam in quarta parte secundi incipit. Omnia sunt in capsulis
usuarii:

\begin{longtable}{@{}llp{0.6\textwidth}@{}}
\toprule
Quid & Systema & Ubi\\
\midrule
\endhead
Optiones & Linux & \texttt{\textasciitilde/.config/py-words/config.toml}\\
 & macOS & \texttt{\textasciitilde/Library/Application Support/py-words}\\
 & Windows & \texttt{\%APPDATA\%\textbackslash py-words}\\
\addlinespace
Historia & Linux & \texttt{\textasciitilde/.local/state/py-words/history}\\
 & macOS & \texttt{\textasciitilde/Library/Application Support/py-words}\\
 & Windows & \texttt{\%LOCALAPPDATA\%\textbackslash py-words}\\
\addlinespace
Cella & Linux & \texttt{\textasciitilde/.cache/py-words}\\
 & macOS & \texttt{\textasciitilde/Library/Caches/py-words}\\
 & Windows & \texttt{\%LOCALAPPDATA\%\textbackslash py-words}\\
\bottomrule
\end{longtable}

\noindent\index{variabiles ambitus}%
\texttt{\$PY\_WORDS\_CONFIG}, \texttt{\$PY\_WORDS\_HISTORY} et
\texttt{\$PY\_WORDS\_CACHE} alia loca dant; \texttt{\$PY\_WORDS\_HISTORY}
vacua nullam historiam servat, et \texttt{\$PY\_WORDS\_CACHE} vacua nullam
cellam. Cella semper deleri potest. Fasciculi datorum queruntur primo per
\texttt{--datadir}, deinde in \texttt{\$PY\_WORDS\_DATADIR}, deinde in
capsula quam fasciculus configurationis nominat, et ultimo in capsula
\texttt{data} iuxta programma.

\section{Cum aliquid deficit}

\index{remedia}%
\begin{description}
\item[Programma dicit Python esse nimis antiquum] institue Python 3.13
  vel posterius, vel utere AppImage.
\item[Fenestra non aperitur: no module named tkinter] institue Tk
  (capitulum~\ref{sec:install}).
\item[Verbum non invenitur] inspice \menu{Inquisitio approximata}, et optiones
  \texttt{use\_tricks} et \texttt{medieval\_spellings}; verbum ignotum
  cum littera maiuscula habetur pro nomine proprio, nisi
  \texttt{ignore\_unknown\_names} sit ablata.
\item[Versus inventa mira dat] aufer \texttt{trim\_output}, que formas
  insolitas omittit.
\item[Logeion non respondet] interrete requirit; \menu{Quere in\then
  Hic} ad hoc dictionarium redit.
\item[\menu{Emenda dictionarium\ldots} est cinereum] data scribi non possunt
  (AppImage, RPM), vel Logeion electum est (capitulum~\ref{sec:editing}).
\item[Textus fenestre est nimis parvus vel magnus] \key{Ctrl++} et
  \key{Ctrl+-}, vel \menu{Aspectus\then Typi litterarum\ldots}.
\item[Optio male se habet] dele \texttt{config.toml}
  (capitulum~\ref{sec:files}); programma incipit cum optionibus
  primitivis.
\end{description}

\noindent
Difficultates et vota scribi possunt tamquam \textit{issues} in GitHub vel
in Codeberg (loca sunt in prima pagina), vel mitti ad
\href{mailto:nicolasvaughan@proton.me}{\texttt{nicolasvaughan@proton.me}}.

\printindex

\end{document}

% ===============================================
%%% Local Variables:
%%% mode: latex
%%% lsp-ltex-plus-language: "la"
%%% TeX-master: t
%%% End:
